\subsubsection{Bernstein-Vazirani's Algorithm}\label{sec:bernstein-vazirani-algorithm}

\paragraph{Problem Definition and Classical Solution}\label{sec:bernstein-vazirani-algorithm-problem-definition-and-classical-solution}

\definition{Bernstein-Vazirani's Algorithm}, proposed in 1992 by Ethan Bernstein and Umesh Vazirani, is one of the first examples showing that a \hl{quantum computer can solve a problem with dramatically fewer oracle queries than any deterministic classical algorithm.}

\highspace
Like Deutsch's algorithm, it is formulated as an \textbf{oracle problem}. The oracle hides some information, and our goal is to discover it by querying the oracle as few times as possible.

\highspace
\begin{flushleft}
    \textcolor{Green3}{\faIcon{tools} \textbf{The Hidden Bitstring}}
\end{flushleft}
Suppose there exists an unknown binary string:
\begin{equation*}
    w \in \left\{0,1\right\}^{n}
\end{equation*}
Called the \textbf{hidden bitstring}. For example, if $n=5$, a possible hidden bitstring could be:
\begin{equation*}
    w = 10110
\end{equation*}
The algorithm \textbf{knows nothing} about $w$. It can \hl{only obtain information by querying an oracle.}

\highspace
\begin{flushleft}
    \textcolor{Green3}{\faIcon{magic} \textbf{The Oracle Function}}
\end{flushleft}
The oracle implements the Boolean function:
\begin{equation*}
    f_w (x) = w \cdot x \; \left(\mod 2\right)
\end{equation*}
Where:
\begin{itemize}
    \item $x \in \left\{0,1\right\}^{n}$ is an $n$-bit input string chosen by us.
    \item $w$ is the unknown \emph{hidden bitstring}.
    \item $w \cdot x$ denotes the \textbf{bitwise inner product modulo 2}. The inner product module 2 is computed by:
    \begin{equation*}
        w \cdot x = \left(w_1 x_1\right) \oplus \left(w_2 x_2\right) \oplus \cdots \oplus \left(w_n x_n\right)
    \end{equation*}
    Where multiplication is the ordinary logical AND, and $\oplus$ is the logical XOR. For example:
    \begin{equation*}
        w = 101, \quad x = 110
    \end{equation*}
    Then:
    \begin{equation*}
        w \cdot x = \left(1 \cdot 1\right) \oplus \left(0 \cdot 1\right) \oplus \left(1 \cdot 0\right) = 1 \oplus 0 \oplus 0 = 1 \oplus 0 = 1
    \end{equation*}
    Therefore:
    \begin{equation*}
        f_w (110) = 1
    \end{equation*}
    Similarly, if $x = 011$, then:
    \begin{equation*}
        w \cdot x = \left(1 \cdot 0\right) \oplus \left(0 \cdot 1\right) \oplus \left(1 \cdot 1\right) = 0 \oplus 0 \oplus 1 = 0 \oplus 1 = 1
    \end{equation*}
    Therefore:
    \begin{equation*}
        f_w (011) = 1
    \end{equation*}
\end{itemize}

\highspace
\begin{flushleft}
    \textcolor{Green3}{\faIcon{balance-scale} \textbf{Goal of the Algorithm and Classical Deterministic Solution}}
\end{flushleft}
Our objective is \textbf{not} to evaluate the function for one particular input. Instead, the \hl{goal is to determine the entire hidden string $w$ using as few oracle queries as possible.} Once $w$ is known, we can compute $f_w (x)$ for any input $x$ ourselves, without querying the oracle again. This is the only information hidden inside the oracle.

\highspace
Instead, a \hl{deterministic classical algorithm must determine every bit of $w$.} The simplest strategy is to query the oracle with inputs containing a single 1. For example:
\begin{equation*}
    x = 100 \dots 0
\end{equation*}
Then:
\begin{equation*}
    f_w (100 \dots 0) = w_1
\end{equation*}
Likewise, querying the oracle with:
\begin{equation*}
    x = 010 \dots 0, \quad f_w (010 \dots 0) = w_2
\end{equation*}
And so on, until:
\begin{equation*}
    x = 000 \dots 1, \quad f_w (000 \dots 1) = w_n
\end{equation*}
Each query reveals exactly one bit of the hidden string. Therefore, \hl{recovering all $n$ bits requires exactly $n$ oracle queries.} For example:
\begin{itemize}
    \item Query 1: $f_w (1000) = 1$
    \item Query 2: $f_w (0100) = 0$
    \item Query 3: $f_w (0010) = 1$
    \item Query 4: $f_w (0001) = 1$
\end{itemize}
After 4 oracle calls, the algorithm reconstructs the hidden string:
\begin{equation*}
    w = 1011
\end{equation*}
In general, an $n$-bit hidden string requires $n$ queries.

\highspace
\begin{flushleft}
    \textcolor{Green3}{\faIcon{question-circle} \textbf{Why This Problem Is Interesting}}
\end{flushleft}
At first glance, this problem appears impossible to accelerate: every oracle call seems to reveal only one bit of information about the hidden string. Surprisingly, a quantum computer can exploit \textbf{superposition} and \textbf{phase kickback} to extract \textbf{all $n$ bits simultaneously}. The Bernstein-Vazirani algorithm determines the complete hidden bitstring using \textbf{only 1 oracle query}, providing a clear example of a \textbf{query-complexity advantage} over deterministic classical computation.

\highspace
\begin{takeawaysbox}[: Bernstein-Vazirani Problem]
    \textbf{Bernstein-Vazirani problem}: an oracle hides an unknown bitstring $w$ through the function:
    \begin{equation*}
        f_w (x) = w \cdot x \; \left(\mod 2\right)
    \end{equation*}
    The objective is to recover $w$. A deterministic classical algorithm requires $n$ oracle queries (one per bit), whereas the quantum algorithm recovers the entire string with a single oracle query by exploiting superposition and phase kickback.
\end{takeawaysbox}